\section{Einleitung und Problemhintergrund}

Pressemitteilungen der Statistik Austria wie \glqq \acs{pkw}-Bestand 2025 um 54.208 höher als 2024\grqq{} oder \glqq 2025 so viele Pkw-Neuzulassungen wie seit 2019 nicht mehr\grqq{} bestätigen den Trend der zunehmenden Motorisierung innerhalb Österreichs \parencites[S. 1]{statistikaustriaPkwBestand2025Um2026}[S. 1]{statistikaustria2025VielePkwNeuzulassungen2026}.
Mit mehr als 7,5 Millionen zugelassenen \acp{kfz} wurde zwischen 1990 und 2025 ein Höchststand erreicht, wobei besonders große und schwere Fahrzeuge wie Geländewagen und \acp{suv} immer mehr Platz im öffentlichen Raum einnehmen \parencite{hatzlDatenZeigenImmer2025}.
Daraus lässt sich schließen, dass der fließende und stehende Verkehr innerhalb von Städten optimiert werden muss, um dem erhöhten Verkehrsaufkommen entgegenzuwirken \parencite{jaschinskyWarumOsterreichsStadte2025}.

Während Wien selbst einen geringeren \acs{kfz}-Bestand als andere Bundesländer aufweist, pendeln Erwerbstätige aus dem Umland regelmäßig nach Wien \parencite{statistikaustriaKfzBestand2026}.
Aufgrund mangelnder Anbindungen an den öffentlichen Verkehr erfolgt die Reise unter Umständen mit dem privaten \acs{kfz} \parencite{aknoAKNiederosterreichPendleranalyse20252025}.
Dabei endet die Suche nach einem geeigneten Parkplatz unter anderem in verschiedenen Parkhäusern, die nur teilweise mit Parkleitsystemen ausgestattet sind. Diese Systeme sollen vor Ort bei der Suche nach einem geeignetem Parkplatz helfen und geben beispielsweise Auskunft über die Anzahl der freien Stellplätze innerhalb des Parkhauses. Speziell in Wien dienen Funktionszonen als Orientierungshilfe, in denen darin befindliche Parkäuser als Einheit erfasst und Parkleitsysteme durch Leitfarben einheitlich gestaltet werden \parencite{stadtwienParkleitsystemFuerWien2026}.

Vor Fahrtantritt wird die Reise stellenweise durch Google Maps und andere Kartendienste geplant, wobei dort keine aktuellen Informationen über die Parksituation abgerufen werden können. Stattdessen ist das Aufrufen unterschiedlicher Websites notwendig, auf denen die gewünschte Information erst gefunden werden muss. Bei akuten Änderungen der Parksituation (z.B. das Parkhaus wird während der Anfahrt vollständig besetzt) wäre es von Vorteil, die wichtigsten Informationen auch während der Zielführung abrufen zu können.
Hinsichtlich des Ablenkungspotentials durch Smartphones ist es während der Fahrt meist nicht problemlos möglich, sich mit solchen akuten Änderungen auseinanderzusetzen \parencite[S. 361]{oviedo-trespalaciosUnderstandingImpactsMobile2016}.
In Österreich zeigt die Verkehrsstatistik aus dem Jahr 2025, dass von nahezu 400 Verkehrstoten die Unfallursachen zu knapp einem Drittel auf Unachtsamkeit und Ablenkung zurückzuführen sind \parencite{bmi397VerkehrstoteAuf2026}.
Navigieren, lesen und schreiben zählen zu den namhaftesten Interaktionen mit dem Smartphone, die während dem Steuern eines Fahrzeuges einen Einfluss auf Geschwindigkeit, Beschleunigung, Spurhaltung und den Sicherheitsabstand haben können. Auch die Leistungsfähigkeit, Informationen aufzunehmen und zu verarbeiten, nimmt beim Fahren und der parallelen Nutzung eines Smartphones ab \parencite[S. 372-373]{oviedo-trespalaciosUnderstandingImpactsMobile2016}.
Die Verwendung adaptiver oder vereinfachter Benutzeroberflächen kann dazu beitragen, Aufgaben schneller zu erledigen und einen besseren Fokus auf den Verkehr zu behalten, ohne die Verwendbarkeit des Smartphones zu sehr einzuschränken \parencite[S. 1-3]{khanDesignContextAwareAdaptive2020}.

Um die notwendigen Interaktionen mit dem Smartphone während der Fahrt auf ein Minimum zu reduzieren, könnte die Entscheidungsfindung selbst erleichtert werden. Alternativen würden damit auf Basis von Empfehlungen und eigenen Präferenzen schlicht bestätigt oder abgelehnt werden. Mit dieser Arbeit soll daher ein erster Lösungsansatz entwickelt werden, um aktuelle Informationen auch während der Zielführung abrufen zu können. Ziel ist es, eine dynamische Entscheidungsmöglichkeit auch während der Wegführung zu bieten. Mithilfe von Sensoren und der Integration in eine \ac{AWS}-Infrastruktur wird der technische Grundstein für ein erweiterbares \ac{IoT}-System gelegt.